\subsection{Discriminative Power of FSM-Based Evaluation (RQ2)}
Having shown that FSM-based evaluation aligns with human judgment on interaction quality, we examine whether it can further discriminate between models and content types, a prerequisite for meaningful comparative analysis.

FSM-derived metrics clearly separate models by capability (Figure~\ref{fig:distribution}). Higher-capacity models achieve consistently stronger scores across structural, semantic, and behavioral dimensions. For example, GPT-5 Mini attains the highest overall score (78.85\%), with notably stronger behavioral coherence (33.13\%), while smaller open-source models such as Qwen-1.5-0.5B score substantially lower overall (49.33\%) and on behavioral coherence (10.29\%). Mid-tier models, including GPT-3.5-Turbo and DeepSeek-V3, cluster closely in overall score (70.31\% vs. 70.13\%) but differ in structural and behavioral profiles, indicating nuanced differences in interaction stability rather than simple code correctness.

To quantify the statistical significance of performance differences between adjacent model tiers, we conducted pairwise Welch's \textit{t}-tests on FSM similarity scores. Welch's test is appropriate for comparing means between two groups with potentially unequal variances and sample sizes~\cite{welch1947}. For each pair of consecutive models in the ranking (e.g., GPT-5 Mini vs. GPT-3.5-Turbo), we compute the \textit{t}-statistic as:
\begin{equation}
t = \frac{\bar{x}_1 - \bar{x}_2}{\sqrt{\frac{s_1^2}{n_1} + \frac{s_2^2}{n_2}}}
\end{equation}
where $\bar{x}_i$, $s_i^2$, and $n_i$ represent the mean, variance, and sample size of group $i$, respectively. The Welch-Satterthwaite equation approximates the degrees of freedom. The resulting \textit{p}-value indicates the probability of observing such a difference under the null hypothesis of equal means; $p < 0.05$ is conventionally considered statistically significant.

Results reveal highly significant distinctions across major capability tiers (Table~\ref{tab:pairwise-tests}). GPT-5 Mini significantly outperforms GPT-3.5-Turbo ($\Delta\mu = 8.54$, $t(214) = 8.31$, $p < 0.0001$***), demonstrating that frontier models produce measurably superior interactive quality. Conversely, GPT-3.5-Turbo and DeepSeek-V3 show no significant difference ($\Delta\mu = 0.18$, $t(243) = 0.15$, $p = 0.884$), confirming their functional equivalence despite architectural differences. Subsequent tiers---DeepSeek-V3 vs. GPT-4o-Mini ($p < 0.0001$***), GPT-4o-Mini vs. Llama-3.2-1B ($p < 0.0001$***), and Llama-3.2-1B vs. Qwen-1.5-0.5B ($p < 0.0001$***)---all exhibit highly significant separations, validating FSM-based evaluation's discriminative power across the model capability spectrum.

\begin{table}[t]
\centering
\caption{Pairwise \textit{t}-tests between adjacent models. Significance levels: *** $p<0.001$, ** $p<0.01$, * $p<0.05$, n.s.~not significant.}
\label{tab:pairwise-tests}
\small
\begin{tabular}{lrrrrl}
\toprule
\textbf{Model Pair} & \textbf{$\Delta\mu$} & \textbf{$t$} & \textbf{$df$} & \textbf{$p$} & \textbf{Sig.} \\
\midrule
GPT-5-Mini vs GPT-3.5-Turbo & 8.54 & 8.31 & 214 & <0.0001 & *** \\
GPT-3.5-Turbo vs DeepSeek-V3 & 0.18 & 0.15 & 243 & 0.8836 & n.s. \\
DeepSeek-V3 vs GPT-4o-Mini & 6.90 & 5.68 & 241 & <0.0001 & *** \\
GPT-4o-Mini vs Llama-3.2-1B & 6.07 & 4.57 & 226 & <0.0001 & *** \\
Llama-3.2-1B vs Qwen-1.5-0.5B & 7.82 & 4.93 & 165 & <0.0001 & *** \\
\bottomrule
\end{tabular}
\end{table}

FSM evaluation also reveals content-dependent interaction demands. More structured topics, such as data structures and sorting algorithms, yield higher FSM coverage and more stable transition patterns, whereas advanced or procedural topics exhibit greater variance across states and transitions. This suggests that FSM-based evaluation captures task-specific interaction complexity, enabling analysis of how both model capability and content structure shape interactive quality, rather than collapsing performance into a single aggregate score.

These findings establish that FSM-based evaluation possesses strong discriminative validity: it reliably distinguishes models across capability tiers with statistical rigor, identifies functional equivalence where performance converges (e.g., GPT-3.5-Turbo and DeepSeek-V3), and exposes content-specific interaction challenges that aggregate metrics would obscure.
